% !TEX TS-program = xelatex
% Таблица результатов партий и график побед и поражений чёрных.

\documentclass[12pt,a4paper]{article}

\usepackage[margin=1.5cm]{geometry}
\usepackage{fontspec}
\usepackage{polyglossia}
\setdefaultlanguage{russian}
\usepackage[table]{xcolor}
\RequirePackage{tikz}
\usetikzlibrary{arrows.meta}

\IfFontExistsTF{Liberation Serif}{\setmainfont{Liberation Serif}\setsansfont{Liberation Sans}}{%
\IfFontExistsTF{DejaVu Serif}{\setmainfont{DejaVu Serif}\setsansfont{DejaVu Sans}}{%
  \setmainfont{CMU Serif}\setsansfont{CMU Sans Serif}}}
\IfFontExistsTF{Liberation Mono}{\setmonofont{Liberation Mono}}{
\IfFontExistsTF{DejaVu Sans Mono}{\setmonofont{DejaVu Sans Mono}}{%
\setmonofont{CMU Typewriter Text}}}

\setlength{\parindent}{0pt}

% Палитра: нужны цвета (ink и др.) и график на tikz.
% Общая палитра цветов для всего комплекта.
% Меняйте значения HEX и комментарии, имена менять не рекомендуется.
% Стрелки arrow1…arrow4 — условные названия, не переименовывать.

\RequirePackage{xcolor}

% Основная палитра для печати игровой доски и расстановки пешек
\definecolor{ink}{HTML}{26262A}       % текст, рамка доски, подписи клеток, оси в rating
\definecolor{pawn}{HTML}{000000}      % цвет пешек чёрных и контура пешек белых
\definecolor{gridline}{HTML}{B4B4BC}  % внутренние линии сетки
% Отдельный светлый цвет для стрелок на иллюстрациях: такая стрелка должна
% быть светлее чёрной пешки, иначе при взятии (например a2xb3) кончик
% стрелки пропадает на пешке. Поэтому у moveline свой цвет, а не ink.
\definecolor{moveink}{HTML}{B4B4BC}   % стрелки ходов в иллюстрациях и примерах

% Цвета стрелок: каждой стрелке на карточке соответствует конфета того же цвета.
\definecolor{arrow1}{HTML}{00B050}    % зелёный
\definecolor{arrow2}{HTML}{7030A0}    % фиолетовый
\definecolor{arrow3}{HTML}{FF0000}    % красный
\definecolor{arrow4}{HTML}{FFFF00}    % жёлтый


% Колонтитул с лицензией CC BY-SA 4.0 — общий для всех документов.
% Колонтитул с указанием лицензии.
%
% Зачем: CC BY-SA 4.0, раздел 3(a), требует, чтобы сам материал отдавал лицензию и ссылку на неё.
% Комплект печатают и распространяют без README, поэтому указание лицензии должно быть на каждой странице.
%
% Почему без иконки CC:
% официальный значок CC просят не перерисовывать,
% а готовый пакет ccicons лежит в коллекции fontsextra — в CI это лишние полтора гигабайта ради одного значка.
% Текстовое обозначение лицензии со ссылками требование 3(a) выполняет полностью.
%
% Почему без имени:
% комплект перепечатывают и проводят разные организаторы, авторство остаётся в README и на странице релизов.
%
% Подключение:
% % Колонтитул с указанием лицензии.
%
% Зачем: CC BY-SA 4.0, раздел 3(a), требует, чтобы сам материал отдавал лицензию и ссылку на неё.
% Комплект печатают и распространяют без README, поэтому указание лицензии должно быть на каждой странице.
%
% Почему без иконки CC:
% официальный значок CC просят не перерисовывать,
% а готовый пакет ccicons лежит в коллекции fontsextra — в CI это лишние полтора гигабайта ради одного значка.
% Текстовое обозначение лицензии со ссылками требование 3(a) выполняет полностью.
%
% Почему без имени:
% комплект перепечатывают и проводят разные организаторы, авторство остаётся в README и на странице релизов.
%
% Подключение:
% % Колонтитул с указанием лицензии.
%
% Зачем: CC BY-SA 4.0, раздел 3(a), требует, чтобы сам материал отдавал лицензию и ссылку на неё.
% Комплект печатают и распространяют без README, поэтому указание лицензии должно быть на каждой странице.
%
% Почему без иконки CC:
% официальный значок CC просят не перерисовывать,
% а готовый пакет ccicons лежит в коллекции fontsextra — в CI это лишние полтора гигабайта ради одного значка.
% Текстовое обозначение лицензии со ссылками требование 3(a) выполняет полностью.
%
% Почему без имени:
% комплект перепечатывают и проводят разные организаторы, авторство остаётся в README и на странице релизов.
%
% Подключение:
% \input{colophon.tex} в преамбуле,
% затем \pagestyle{ccfoot} (без номеров страниц) или \pagestyle{ccfootpage} (с номерами).

\RequirePackage{xcolor}
\RequirePackage{url}


% =====================================================================================================================
% Единая строка колонтитула: лицензия, ссылка на неё, ссылка на исходники.
%
% Шрифт здесь основной, а не \sffamily: guide.tex и rating.tex не задают \setsansfont,
% и polyglossia падает на латинском Modern Sans без кириллицы.
% Основной шрифт определён во всех документах.
%
% Ссылки написаны прямо в \url, а не через макросы:
% пакет url читает аргумент буквально и \url{\ccrepo} напечатал бы «\ccrepo».
\newcommand{\ccline}{%
  \hbox to \textwidth{%
    \color{black!45}\fontsize{7}{8.5}\selectfont
    \urlstyle{same}%
    CC BY-SA 4.0\enspace\url{https://creativecommons.org/licenses/by-sa/4.0/}%
    \hfil\url{https://github.com/askras/hexapawn-sama/}}%
  \vspace*{1pt}%
}

% =====================================================================================================================
% Тот же колонтитул с номером страницы справа.
% Под номер зарезервировано 2em, иначе \llap наезжает на ссылку на исходники.
\newcommand{\cclinepage}{%
  \hbox to \dimexpr\textwidth-2em\relax{%
    \color{black!45}\fontsize{7}{8.5}\selectfont
    \urlstyle{same}%
    CC BY-SA 4.0\enspace\url{https://creativecommons.org/licenses/by-sa/4.0/}%
    \hfil\url{https://github.com/askras/hexapawn-sama/}}%
  \hfil\llap{\color{black!45}\fontsize{7}{8.5}\selectfont\thepage}%
  \vspace*{1pt}%
}

% =====================================================================================================================
% Поля узкие (1,2 см), поэтому колонтитул ставим ближе к тексту, иначе он уходит за край листа.
\setlength{\footskip}{9mm}

% =====================================================================================================================
\makeatletter
% Вариант без номеров страниц: карточки, поле, таблица результатов.
\def\ps@ccfoot{%
  \let\@oddhead\@empty
  \let\@evenhead\@empty
  \def\@oddfoot{\ccline}%
  \def\@evenfoot{\ccline}%
}
% Вариант с номерами: методичка, правила, примеры.
\def\ps@ccfootpage{%
  \let\@oddhead\@empty
  \let\@evenhead\@empty
  \def\@oddfoot{\cclinepage}%
  \def\@evenfoot{\cclinepage}%
}
\makeatother
 в преамбуле,
% затем \pagestyle{ccfoot} (без номеров страниц) или \pagestyle{ccfootpage} (с номерами).

\RequirePackage{xcolor}
\RequirePackage{url}


% =====================================================================================================================
% Единая строка колонтитула: лицензия, ссылка на неё, ссылка на исходники.
%
% Шрифт здесь основной, а не \sffamily: guide.tex и rating.tex не задают \setsansfont,
% и polyglossia падает на латинском Modern Sans без кириллицы.
% Основной шрифт определён во всех документах.
%
% Ссылки написаны прямо в \url, а не через макросы:
% пакет url читает аргумент буквально и \url{\ccrepo} напечатал бы «\ccrepo».
\newcommand{\ccline}{%
  \hbox to \textwidth{%
    \color{black!45}\fontsize{7}{8.5}\selectfont
    \urlstyle{same}%
    CC BY-SA 4.0\enspace\url{https://creativecommons.org/licenses/by-sa/4.0/}%
    \hfil\url{https://github.com/askras/hexapawn-sama/}}%
  \vspace*{1pt}%
}

% =====================================================================================================================
% Тот же колонтитул с номером страницы справа.
% Под номер зарезервировано 2em, иначе \llap наезжает на ссылку на исходники.
\newcommand{\cclinepage}{%
  \hbox to \dimexpr\textwidth-2em\relax{%
    \color{black!45}\fontsize{7}{8.5}\selectfont
    \urlstyle{same}%
    CC BY-SA 4.0\enspace\url{https://creativecommons.org/licenses/by-sa/4.0/}%
    \hfil\url{https://github.com/askras/hexapawn-sama/}}%
  \hfil\llap{\color{black!45}\fontsize{7}{8.5}\selectfont\thepage}%
  \vspace*{1pt}%
}

% =====================================================================================================================
% Поля узкие (1,2 см), поэтому колонтитул ставим ближе к тексту, иначе он уходит за край листа.
\setlength{\footskip}{9mm}

% =====================================================================================================================
\makeatletter
% Вариант без номеров страниц: карточки, поле, таблица результатов.
\def\ps@ccfoot{%
  \let\@oddhead\@empty
  \let\@evenhead\@empty
  \def\@oddfoot{\ccline}%
  \def\@evenfoot{\ccline}%
}
% Вариант с номерами: методичка, правила, примеры.
\def\ps@ccfootpage{%
  \let\@oddhead\@empty
  \let\@evenhead\@empty
  \def\@oddfoot{\cclinepage}%
  \def\@evenfoot{\cclinepage}%
}
\makeatother
 в преамбуле,
% затем \pagestyle{ccfoot} (без номеров страниц) или \pagestyle{ccfootpage} (с номерами).

\RequirePackage{xcolor}
\RequirePackage{url}


% =====================================================================================================================
% Единая строка колонтитула: лицензия, ссылка на неё, ссылка на исходники.
%
% Шрифт здесь основной, а не \sffamily: guide.tex и rating.tex не задают \setsansfont,
% и polyglossia падает на латинском Modern Sans без кириллицы.
% Основной шрифт определён во всех документах.
%
% Ссылки написаны прямо в \url, а не через макросы:
% пакет url читает аргумент буквально и \url{\ccrepo} напечатал бы «\ccrepo».
\newcommand{\ccline}{%
  \hbox to \textwidth{%
    \color{black!45}\fontsize{7}{8.5}\selectfont
    \urlstyle{same}%
    CC BY-SA 4.0\enspace\url{https://creativecommons.org/licenses/by-sa/4.0/}%
    \hfil\url{https://github.com/askras/hexapawn-sama/}}%
  \vspace*{1pt}%
}

% =====================================================================================================================
% Тот же колонтитул с номером страницы справа.
% Под номер зарезервировано 2em, иначе \llap наезжает на ссылку на исходники.
\newcommand{\cclinepage}{%
  \hbox to \dimexpr\textwidth-2em\relax{%
    \color{black!45}\fontsize{7}{8.5}\selectfont
    \urlstyle{same}%
    CC BY-SA 4.0\enspace\url{https://creativecommons.org/licenses/by-sa/4.0/}%
    \hfil\url{https://github.com/askras/hexapawn-sama/}}%
  \hfil\llap{\color{black!45}\fontsize{7}{8.5}\selectfont\thepage}%
  \vspace*{1pt}%
}

% =====================================================================================================================
% Поля узкие (1,2 см), поэтому колонтитул ставим ближе к тексту, иначе он уходит за край листа.
\setlength{\footskip}{9mm}

% =====================================================================================================================
\makeatletter
% Вариант без номеров страниц: карточки, поле, таблица результатов.
\def\ps@ccfoot{%
  \let\@oddhead\@empty
  \let\@evenhead\@empty
  \def\@oddfoot{\ccline}%
  \def\@evenfoot{\ccline}%
}
% Вариант с номерами: методичка, правила, примеры.
\def\ps@ccfootpage{%
  \let\@oddhead\@empty
  \let\@evenhead\@empty
  \def\@oddfoot{\cclinepage}%
  \def\@evenfoot{\cclinepage}%
}
\makeatother

\pagestyle{ccfoot}

\begin{document}


\begin{center}
{\LARGE\bfseries Таблица результатов партий}\\[0.2em]
{\small Поставьте галочку в столбец победителя каждой партии}\\[0.6em]

\renewcommand{\arraystretch}{0.95}
\footnotesize
\begin{minipage}[t]{0.47\textwidth}
\begin{tabular}{|>{\centering\arraybackslash}m{1.3cm}|>{\centering\arraybackslash}m{0.33\textwidth}|>{\centering\arraybackslash}m{0.33\textwidth}|}
\hline
\rowcolor{black!92}\textbf{\textcolor{white}{\parbox{1.3cm}{\centering Номер\\[-1pt]партии}}} & \textbf{\textcolor{white}{\parbox{0.33\textwidth}{\centering Белые\\[-1pt](гроссмейстер)}}} & \textbf{\textcolor{white}{\parbox{0.33\textwidth}{\centering Чёрные\\[-1pt](машина)}}}\\
\hline
1 & & \\
\hline
2 & & \\
\hline
3 & & \\
\hline
4 & & \\
\hline
5 & & \\
\hline
6 & & \\
\hline
7 & & \\
\hline
8 & & \\
\hline
9 & & \\
\hline
10 & & \\
\hline
11 & & \\
\hline
12 & & \\
\hline
13 & & \\
\hline
14 & & \\
\hline
15 & & \\
\hline
16 & & \\
\hline
17 & & \\
\hline
18 & & \\
\hline
19 & & \\
\hline
20 & & \\
\hline
21 & & \\
\hline
22 & & \\
\hline
23 & & \\
\hline
24 & & \\
\hline
25 & & \\
\hline
26 & & \\
\hline
27 & & \\
\hline
28 & & \\
\hline
29 & & \\
\hline
30 & & \\
\hline
31 & & \\
\hline
32 & & \\
\hline
33 & & \\
\hline
34 & & \\
\hline
35 & & \\
\hline
36 & & \\
\hline
37 & & \\
\hline
38 & & \\
\hline
39 & & \\
\hline
40 & & \\
\hline
41 & & \\
\hline
42 & & \\
\hline
43 & & \\
\hline
44 & & \\
\hline
45 & & \\
\hline
46 & & \\
\hline
47 & & \\
\hline
48 & & \\
\hline
49 & & \\
\hline
50 & & \\
\hline

\end{tabular}
\end{minipage}\hfill
\begin{minipage}[t]{0.47\textwidth}
\begin{tabular}{|>{\centering\arraybackslash}m{1.3cm}|>{\centering\arraybackslash}m{0.33\textwidth}|>{\centering\arraybackslash}m{0.33\textwidth}|}
\hline
\rowcolor{black!92}\textbf{\textcolor{white}{\parbox{1.3cm}{\centering Номер\\[-1pt]партии}}} & \textbf{\textcolor{white}{\parbox{0.33\textwidth}{\centering Белые\\[-1pt](гроссмейстер)}}} & \textbf{\textcolor{white}{\parbox{0.33\textwidth}{\centering Чёрные\\[-1pt](машина)}}}\\
\hline
51 & & \\
\hline
52 & & \\
\hline
53 & & \\
\hline
54 & & \\
\hline
55 & & \\
\hline
56 & & \\
\hline
57 & & \\
\hline
58 & & \\
\hline
59 & & \\
\hline
60 & & \\
\hline
61 & & \\
\hline
62 & & \\
\hline
63 & & \\
\hline
64 & & \\
\hline
65 & & \\
\hline
66 & & \\
\hline
67 & & \\
\hline
68 & & \\
\hline
69 & & \\
\hline
70 & & \\
\hline
71 & & \\
\hline
72 & & \\
\hline
73 & & \\
\hline
74 & & \\
\hline
75 & & \\
\hline
76 & & \\
\hline
77 & & \\
\hline
78 & & \\
\hline
79 & & \\
\hline
80 & & \\
\hline
81 & & \\
\hline
82 & & \\
\hline
83 & & \\
\hline
84 & & \\
\hline
85 & & \\
\hline
86 & & \\
\hline
87 & & \\
\hline
88 & & \\
\hline
89 & & \\
\hline
90 & & \\
\hline
91 & & \\
\hline
92 & & \\
\hline
93 & & \\
\hline
94 & & \\
\hline
95 & & \\
\hline
96 & & \\
\hline
97 & & \\
\hline
98 & & \\
\hline
99 & & \\
\hline
100 & & \\
\hline

\end{tabular}
\end{minipage}
\end{center}

\vfill

\begin{center}
{\small \emph{Подсказка:}
после первых 50 партий постройте график побед и поражений чёрных (машины) и посмотрите, как меняется число выигранных
машин партий по мере обучения.}
\end{center}

\newpage

% сетка: X 0..100 (партии), Y -50..+50 (счёт = победы чёрных минус победы белых)
\newcommand{\maxX}{100}
\newcommand{\maxY}{50}

\vspace*{\fill}

\begin{center}
{\LARGE\bfseries График побед и поражений чёрных}\\[0.2em]
{\small ось X --- номер партии, ось Y --- счёт: победы чёрных минус победы белых}
\end{center}

\begin{center}
\begin{tikzpicture}[x=0.155cm, y=0.20cm, line width=0.4pt]
  % мелкая сетка: каждое деление
  \foreach \x in {0,...,100} {\draw[black!12] (\x, -\maxY) -- (\x, \maxY);}
  \foreach \y in {-\maxY,...,0} {\draw[black!12] (0, \y) -- (\maxX, \y);}
  \foreach \y in {1,...,\maxY} {\draw[black!12] (0, \y) -- (\maxX, \y);}
  % средняя (нулевая) линия с засечками через 5 (без подписей)
  \draw[black!55] (0, 0) -- (\maxX, 0);
  \foreach \x in {5,10,15,...,95} {\draw[black!55] (\x, -0.5) -- (\x, 0.5);}

  % рамка поля
  \draw[ink, line width=0.8pt] (0, -\maxY) rectangle (\maxX, \maxY);

  % оси
  \draw[-{Stealth[length=2.5mm,width=1.8mm]}, ink, line width=0.8pt]
    (-2, -\maxY) -- (\maxX+3, -\maxY) node[below, pos=1, xshift=0.1cm] {};
  \draw[-{Stealth[length=2.5mm,width=1.8mm]}, ink, line width=0.8pt]
    (0, -\maxY-2) -- (0, \maxY+3);

  % подписи осей
  \node[below, font=\scriptsize] at (\maxX/2, -\maxY-4) {номер партии};
  \node[rotate=90, font=\scriptsize] at (-8, 0) {счёт (победы чёрных $-$ победы белых)};
  \node[below left, font=\scriptsize] at (0, -\maxY-2) {0};

  % числовые отметки по Y: подписи слева (-50...50) и справа (-45...50, без -50)
  \foreach \y in {-50,-45,...,50} {
    \draw (0, \y) -- (0.6, \y);
    \node[left, font=\scriptsize, xshift=-0.12cm] at (0, \y) {\y};
  }
  \foreach \y in {-45,-40,...,50} {
    \draw (\maxX, \y) -- ({\maxX-0.6}, \y);
    \node[right, font=\scriptsize, xshift=0.12cm] at (\maxX, \y) {\y};
  }
  % числовые отметки по X: подписи снизу и сверху
  \foreach \x in {5,10,15,...,95} {
    \draw (\x, -\maxY) -- (\x, {-\maxY+0.6});
    \draw (\x, \maxY) -- (\x, {\maxY-0.6});
    \node[below, font=\scriptsize, yshift=-0.12cm] at (\x, -\maxY) {\x};
    \node[above, font=\scriptsize, yshift=0.12cm] at (\x, \maxY) {\x};
  }

  % начальная точка (0; 0) --- от неё строится график
  \fill[black] (0, 0) circle[radius=0.1cm];

  % образец одного хода (вправо +1, вверх +1) --- наклонный отрезок;
  % раскомментируйте, чтобы увидеть стиль линии графика
  % \draw[red!75!black, line width=0.6pt, -] (0, 0) -- (1, 1) -- (2, 0) -- (3, 1);
  % \foreach \px/\py in {0/0, 1/1, 2/0, 3/1} \fill[red!75!black] (\px, \py) circle[radius=0.12cm];
\end{tikzpicture}
\end{center}

\vspace*{\fill}

\begin{center}
{\footnotesize\itshape Как строить график: начиная с точки (0; 0), для каждой партии откладывайте по оси X один шаг
вправо, а по оси Y --- на одно деление вверх (победа чёрных) или на одно деление вниз (победа белых).
Отмечайте точки и соединяйте их наклонными прямыми отрезками (не ступеньками):
каждый отрезок идёт из текущей точки вправо-вверх или вправо-вниз.}
\end{center}

\end{document}
